\textbf{Phần 1. Trắc nghiệm lựa chọn}
\setcounterpageref{ex}{0}

\begin{ex} %[2H3N15]
Ki-lô-gam được viết tắt là
\choice
{$\text{km}$}
{$\text{g}$}
{$\text{l}$}
{\True $\text{kg}$}
\loigiai{
Ki-lô-gam viết tắt là $\text{kg}$.\\
\textbf{Đáp án của câu hỏi là D}
}
\end{ex}

\begin{ex} %[2H3N15]
Để đo khối lượng (cân nặng) của một vật, người ta dùng đơn vị
\choice
{Xăng-ti-mét}
{Lít}
{\True Ki-lô-gam}
{Giờ}
\loigiai{
Đo khối lượng dùng ki-lô-gam.\\
\textbf{Đáp án của câu hỏi là C}
}
\end{ex}

\begin{ex} %[2H3H15]
Tính: $15\text{ kg} + 8\text{ kg} = \dots$
\choice
{$22\text{ kg}$}
{\True $23\text{ kg}$}
{$24\text{ kg}$}
{$25\text{ kg}$}
\loigiai{
$15 + 8 = 23\text{ kg}$.\\
\textbf{Đáp án của câu hỏi là B}
}
\end{ex}

\begin{ex} %[2H3H15]
Tính: $35\text{ kg} - 12\text{ kg} = \dots$
\choice
{\True $23\text{ kg}$}
{$22\text{ kg}$}
{$24\text{ kg}$}
{$20\text{ kg}$}
\loigiai{
$35 - 12 = 23\text{ kg}$.\\
\textbf{Đáp án của câu hỏi là A}
}
\end{ex}

\begin{ex} %[2H3H15]
Bao gạo to nặng $30\text{ kg}$, bao gạo nhỏ nặng $15\text{ kg}$. Cả hai bao gạo nặng số ki-lô-gam là
\choice
{$40\text{ kg}$}
{\True $45\text{ kg}$}
{$50\text{ kg}$}
{$55\text{ kg}$}
\loigiai{
$30 + 15 = 45\text{ kg}$.\\
\textbf{Đáp án của câu hỏi là B}
}
\end{ex}

\begin{ex} %[2H3H15]
Con lợn cân nặng $65\text{ kg}$, con chó cân nặng $15\text{ kg}$. Con lợn nặng hơn con chó số ki-lô-gam là
\choice
{$40\text{ kg}$}
{$60\text{ kg}$}
{\True $50\text{ kg}$}
{$80\text{ kg}$}
\loigiai{
$65 - 15 = 50\text{ kg}$.\\
\textbf{Đáp án của câu hỏi là C}
}
\end{ex}

\begin{ex} %[2H3V15]
Một quả dưa hấu cân nặng $4\text{ kg}$, quả bí đỏ cân nặng $5\text{ kg}$. Cả hai quả cân nặng là
\choice
{$8\text{ kg}$}
{$10\text{ kg}$}
{$11\text{ kg}$}
{\True $9\text{ kg}$}
\loigiai{
$4 + 5 = 9\text{ kg}$.\\
\textbf{Đáp án của câu hỏi là D}
}
\end{ex}

\begin{ex} %[2H3V15]
Có một bao đường nặng $20\text{ kg}$, người ta đã lấy ra $8\text{ kg}$. Số đường còn lại trong bao là
\choice
{$11\text{ kg}$}
{\True $12\text{ kg}$}
{$13\text{ kg}$}
{$14\text{ kg}$}
\loigiai{
$20 - 8 = 12\text{ kg}$.\\
\textbf{Đáp án của câu hỏi là B}
}
\end{ex}

\begin{ex} %[2H3V15]
Điền dấu thích hợp: $18\text{ kg} + 5\text{ kg} \dots 25\text{ kg}$.
\choice
{$>$}
{$=$}
{\True $<$}
{Không so sánh được}
\loigiai{
$18\text{ kg} + 5\text{ kg} = 23\text{ kg} < 25\text{ kg}$.\\
\textbf{Đáp án của câu hỏi là C}
}
\end{ex}

\begin{ex} %[2H3V15]
Tính giá trị biểu thức: $40\text{ kg} + 25\text{ kg} - 15\text{ kg}$.
\choice
{\True $50\text{ kg}$}
{$45\text{ kg}$}
{$55\text{ kg}$}
{$60\text{ kg}$}
\loigiai{
$65 - 15 = 50\text{ kg}$.\\
\textbf{Đáp án của câu hỏi là A}
}
\end{ex}

\begin{ex} %[2H3V15]
Mẹ mua về $15\text{ kg}$ gạo tẻ và $8\text{ kg}$ gạo nếp. Hỏi mẹ mua tất cả bao nhiêu ki-lô-gam gạo?
\choice
{$22\text{ kg}$}
{$24\text{ kg}$}
{$25\text{ kg}$}
{\True $23\text{ kg}$}
\loigiai{
$15 + 8 = 23\text{ kg}$.\\
\textbf{Đáp án của câu hỏi là D}
}
\end{ex}

\begin{ex} %[2H3C15]
Có một cân đĩa và các quả cân loại $1\text{ kg}, 2\text{ kg}, 5\text{ kg}$. Để cân một túi hoa quả nặng $8\text{ kg}$, ta có thể đặt ở đĩa bên kia các quả cân nào?
\choice
{\True $5\text{ kg}, 2\text{ kg}$ và $1\text{ kg}$}
{$5\text{ kg}$ và $2\text{ kg}$}
{2 quả $5\text{ kg}$}
{3 quả $2\text{ kg}$}
\loigiai{
$5 + 2 + 1 = 8\text{ kg}$.\\
\textbf{Đáp án của câu hỏi là A}
}
\end{ex}

\textbf{Phần 2. Trắc nghiệm đúng sai}
\setcounterpageref{ex}{0}

\begin{ex} %[2H3H15]
Xét tính đúng sai của các phép tính với đơn vị đo:
\choiceTF
{\True $12\text{ kg} + 8\text{ kg} = 20\text{ kg}$.}
{\True $30\text{ kg} - 10\text{ kg} = 20\text{ kg}$.}
{$15\text{ kg} + 7 = 22\text{ kg}$.}
{\True Khi tính toán với đại lượng, các số hạng phải cùng kèm đơn vị đo.}
\loigiai{
\begin{itemize}
    \item $12\text{ kg} + 8\text{ kg} = 20\text{ kg}$: Đúng.
    \item $30\text{ kg} - 10\text{ kg} = 20\text{ kg}$: Đúng.
    \item Ý 3 thiếu đơn vị đo ở số thứ hai. Sai.
    \item Cần cùng kèm đơn vị đo: Đúng.
\end{itemize}
\begin{itemize}
    \item a) \textbf{Đúng}.
    \item b) \textbf{Đúng}.
    \item c) \textbf{Sai}.
    \item d) \textbf{Đúng}.
\end{itemize}
}
\end{ex}

\begin{ex} %[2H3V15]
Bao ngô nặng $45\text{ kg}$, bao thóc nặng hơn bao ngô $10\text{ kg}$. Xét tính đúng sai:
\choiceTF
{\True Bao thóc cân nặng $55\text{ kg}$ ($45 + 10 = 55$).}
{\True Cả hai bao cân nặng $100\text{ kg}$ ($45 + 55 = 100$).}
{Bao thóc nhẹ hơn bao ngô.}
{\True Bao ngô nhẹ hơn bao thóc 10 kg.}
\loigiai{
\begin{itemize}
    \item Thóc: $45 + 10 = 55\text{ kg}$: Đúng.
    \item Tổng: $45 + 55 = 100\text{ kg}$: Đúng.
    \item Thóc nặng hơn (không phải nhẹ hơn). Ý 3 Sai, Ý 4 Đúng.
\end{itemize}
\begin{itemize}
    \item a) \textbf{Đúng}.
    \item b) \textbf{Đúng}.
    \item c) \textbf{Sai}.
    \item d) \textbf{Đúng}.
\end{itemize}
}
\end{ex}

\textbf{Phần 3. Trả lời ngắn}
\setcounterpageref{ex}{0}

\begin{ex} %[2H3H15]
Tính: $18\text{ kg} + 9\text{ kg}$ theo ki-lô-gam.
\shortans{27}
\loigiai{
\textbf{Đáp số:} $27$\\
$18 + 9 = 27\text{ kg}$.
}
\end{ex}

\begin{ex} %[2H3H15]
Tính: $45\text{ kg} - 20\text{ kg}$ theo ki-lô-gam.
\shortans{25}
\loigiai{
\textbf{Đáp số:} $25$\\
$45 - 20 = 25\text{ kg}$.
}
\end{ex}

\begin{ex} %[2H3V15]
Bao gạo nặng $35\text{ kg}$, mẹ mua thêm $15\text{ kg}$ nữa. Hỏi có tất cả bao nhiêu ki-lô-gam gạo?
\shortans{50}
\loigiai{
\textbf{Đáp số:} $50$\\
$35 + 15 = 50\text{ kg}$.
}
\end{ex}

\begin{ex} %[2H3V15]
Một con dê nặng $28\text{ kg}$, con cừu nặng $22\text{ kg}$. Hỏi cả hai con nặng bao nhiêu ki-lô-gam?
\shortans{50}
\loigiai{
\textbf{Đáp số:} $50$\\
$28 + 22 = 50\text{ kg}$.
}
\end{ex}

\begin{ex} %[2H3V15]
Một bao xi măng nặng $50\text{ kg}$, thợ xây đã dùng hết $35\text{ kg}$. Hỏi còn lại bao nhiêu ki-lô-gam xi măng?
\shortans{15}
\loigiai{
\textbf{Đáp số:} $15$\\
$50 - 35 = 15\text{ kg}$.
}
\end{ex}

\begin{ex} %[2H3C15]
Quả bưởi nặng $2\text{ kg}$, quả dưa hấu nặng hơn quả bưởi $3\text{ kg}$. Hỏi cả hai quả nặng bao nhiêu ki-lô-gam?
\shortans{7}
\loigiai{
\textbf{Đáp số:} $7$\\
Dưa hấu: $2 + 3 = 5\text{ kg}$. Cả hai quả: $2 + 5 = 7\text{ kg}$.
}
\end{ex}
